% Generated draft: requires mathematical and source review.
\documentclass[11pt]{article}
\usepackage[margin=25mm]{geometry}
\usepackage{fontspec}
\setmainfont{Latin Modern Roman}
\usepackage{amsmath,amssymb,mathtools,unicode-math}
\setmathfont{Latin Modern Math}
\usepackage{enumitem,needspace,xurl,hyperref}
\setlength{\parindent}{0pt}
\setlength{\parskip}{5pt}
\setlength{\emergencystretch}{3em}
\newcommand{\N}{\mathbb{N}}
\newcommand{\Z}{\mathbb{Z}}
\newcommand{\Q}{\mathbb{Q}}
\newcommand{\R}{\mathbb{R}}
\newcommand{\C}{\mathbb{C}}
\newcommand{\F}{\mathbb{F}}
\newcommand{\A}{\mathbb{A}}
\newcommand{\PP}{\mathbb{P}}
\newcommand{\E}{\mathbb{E}}
\newcommand{\eps}{\varepsilon}
\newcommand{\dd}{\,\mathrm{d}}
\begin{document}
\section*{UnsolvedMath ID: 20002881}
\textbf{Source problem code:} catalogue entry
\subsection{Definitions and mathematical statement}
Let $M$ be a strictly pseudoconvex abstract CR manifold of hypersurface type and CR dimension $m$, so $\dim_\R M=2m+1$. In the cooriented setting, let $\pi:C(M)\to M$ be its Fefferman circle bundle and choose a Fefferman metric $g_F$ from its conformal class. Write $\Omega^F$ for the curvature of its Levi--Civita connection and $p_j(\Omega^F)\in\Omega^{4j}(C(M))$ for its Chern--Weil Pontryagin forms. S. Dragomir asks which of these forms obstruct global CR embeddability of $M$. The original question specifies neither a target dimension nor whether the obstruction concerns the differential form itself or its cohomology class.

The supplied aid gives the following cohomological screen. All classes below are taken in real cohomology, as detected by differential forms. The real bundle splitting
\[
TC(M)\simeq\pi^*TM\oplus\R
\]
gives
\[
[p_j(\Omega^F)]=p_j(TC(M))=\pi^*p_j(TM).
\]
If $e\in H^2(M;\R)$ is the real Euler class of the Fefferman circle bundle, the Gysin sequence gives
\[
[p_j(\Omega^F)]=0
\quad\Longleftrightarrow\quad
p_j(TM)\in\operatorname{im}\left(e\smile-:H^{4j-2}(M;\R)\longrightarrow H^{4j}(M;\R)\right).
\]
The contact splitting of $TM$ also gives
\[
p_j(TM)=(-1)^j c_{2j}\bigl(T^{1,0}M\oplus T^{0,1}M\bigr)
\]
in real cohomology. These formulas concern characteristic classes, independently of the metric representative; they do not assert that an exact Pontryagin form vanishes pointwise.

If $M$ admits a realization as a cooriented real hypersurface in $\C^{m+1}$, then $TM\oplus\R$ is trivial. Hence all the displayed classes vanish, and any nonzero $[p_j(\Omega^F)]$ obstructs such a hypersurface realization. This target is narrower than embeddability in arbitrary higher-dimensional $\C^N$; stable triviality is not inferred merely from a CR embedding of higher codimension.

The aid further notes a universal limitation. If $m$ is odd, set $r=(m+1)/2$, so $\dim C(M)=4r$ and $\dim M=4r-1$. Then
\[
[p_r(\Omega^F)]=\pi^*p_r(TM)=0
\]
by degree on the base. In particular, for a CR $3$-manifold ($m=1$), every positive-degree Fefferman Pontryagin class vanishes. Thus these classes cannot distinguish the embeddable and nonembeddable CR structures on $S^3$ mentioned by the aid. This blindness of classes does not decide whether more refined differential-form information supplies obstructions, nor does it classify all obstructions requested by the original question.
\subsection{Short English statement}
Identify Fefferman Pontryagin obstructions to global CR embedding. The supplied aid gives a real-cohomology pullback and Gysin test for obstructions to a cooriented hypersurface realization, while showing that the classes are automatically blind in CR dimension one and in top degree for odd CR dimension.
\subsection{Sources}
\begin{itemize}
\item Local UnsolvedMath registry, permanent ID $20002881$, source code \texttt{catalogue entry}; canonical \texttt{aim-several-complex-variables-notes.json}, \texttt{notes[38]}.
\item AIM workshop \emph{Emerging applications of complexity for CR mappings}, Problem 3, attributed to S. Dragomir: \url{https://aimath.org/WWN/crmappings/crmappings.pdf}. The supplied verification corrects line-break hyphenation in ``pseudocon-vex'' and identifies ``Approximation'' as the next section heading. Statement status is \texttt{corrected\_verified}; original remarks and literature field are empty.
\item The research summary invokes the Fefferman circle-bundle construction and Gysin sequence, and refers to known embeddability distinctions on CR three-spheres, but supplies no authoritative complete bibliographic citations.
\item Research attempt 1: \texttt{partial\_result}, status at run \texttt{open}, current status \texttt{partially\_solved}, classification \texttt{PARTIAL-PROGRESS}, rigor \texttt{proved}, novelty confidence low. Its cohomological screen is machine-generated aid, not peer reviewed, and requires independent verification before citation.
\end{itemize}
\end{document}
